% Conference paper. Every number below is listed in paper/NOMBRES.md with the line of
% docs/results_summary.md it comes from. Build: latexmk (see paper/latexmkrc).
%
% --- Venue parameters: change these lines if the target conference changes ------------
\documentclass[conference]{IEEEtran}   % IEEE conference, two columns
% Page limit targeted: 6 pages including references.
% ----------------------------------------------------------------------------------------

% Reproducible PDF: no creation date, no random trailer ID.
\pdfinfoomitdate=1
\pdftrailerid{}

\usepackage[T1]{fontenc}
\usepackage{amsmath,amssymb}
\usepackage{graphicx}
\usepackage{booktabs}
\usepackage{array}
\usepackage{cite}
\usepackage{url}
\usepackage{tikz}
\usetikzlibrary{arrows.meta}

\newcommand{\figdir}{figures}

\begin{document}

\title{Demand-Matched Exogenous Arrivals Have No Steady State:\\
Lessons from a Port--Steelworks Corridor Rebuilt from Public Data}

\author{\IEEEauthorblockN{Abdessamed Britah}
\IEEEauthorblockA{Independent researcher\\
% TODO(author): affiliation and e-mail address to be confirmed by the author.
}}

\maketitle

\begin{abstract}
Simulation models of port-to-plant supply corridors commonly drive vessel arrivals by an
exogenous process whose mean rate is calibrated on plant demand. We show that such a model
cannot produce a stationary service level: the net input is a random walk without drift,
lost production grows as the square root of the horizon, and every service level it reports
is a transient that depends on the simulated horizon. Storage capacity does not enter the
loss at all. On a corridor rebuilt from public data only, from a port of import to a
direct-reduction steel plant, the measured growth exponent of the cumulative deficit is
0.492 [0.367; 0.599] once the initial stock is removed, and the service level rises from 0.858
at one year to 0.971 at fifty years without levelling off, while multiplying storage capacity by 12 changes production by less
than 0.002\%. The correction is an endogenous supply policy: an $(s,Q)$ order rule on the
inventory position with random lead times. It restores a stationary service level
(0.9867~$\pm$~0.0012 and 0.9912~$\pm$~0.0008 in two scenarios). With discrete furnace
outages, it is not refuted by public traffic statistics, although that test has low power.
We describe how vessel calls were reconstructed from public port snapshots with
interval-censored departures, and illustrate what the corrected model says about investment
levers, without turning it into a recommendation.
\end{abstract}

\begin{IEEEkeywords}
discrete-event simulation, steady state, supply corridor, inventory policy, port operations,
public data
\end{IEEEkeywords}

% =========================================================================================
\section{Introduction}

Bulk import corridors, in which a port receives raw material by sea and forwards it by rail
to a processing plant, are a standard subject for discrete-event simulation
\cite{TODO-port-bulk-simulation,TODO-port-plant-supply-simulation}. A common and natural
modelling choice is to generate vessel arrivals exogenously, at a mean rate equal to plant
demand divided by vessel size, with some dispersion around it, and then to read a service
level, such as the fraction of target production achieved, after one or a few simulated
years \cite{TODO-exogenous-arrival-practice}.

This paper makes one claim. \emph{A model with exogenous arrivals calibrated on demand has no
stationary service level; any value it reports is a transient whose level is set by the
simulated horizon; the remedy is an endogenous supply policy.} The argument is elementary,
but its consequences for practice are large: in such a model, the service level read after
one year depends on the horizon at least as much as on the levers being compared, and
storage capacity, the lever most often studied, has no effect at all.

We found the problem while building a digital twin of an actual corridor from public data
only. The study is independent and not affiliated with the port or the plant operator. The
contributions are:
\begin{enumerate}
\item a method to reconstruct vessel calls from periodic public port snapshots, with
departures treated as interval-censored, and a statement of what such a series can and
cannot identify, complementing reconstructions from vessel-tracking data
\cite{TODO-port-calls-ais} (Section~\ref{sec:reconstruction});
\item a short analytic argument and a measurement showing that demand-matched exogenous
arrivals produce no stationary service level and make storage capacity inert
(Section~\ref{sec:trap});
\item a corrective model with an $(s,Q)$ policy on the inventory position, random lead
times, discrete plant outages and two hypotheses on vessel overtaking, which recovers a
stationary regime and is not refuted by public traffic data
(Sections~\ref{sec:model} and~\ref{sec:results}).
\end{enumerate}

% =========================================================================================
\section{Reconstructing vessel calls from public snapshots}
\label{sec:reconstruction}

\subsection{Data}
The port publishes a situation report as a PDF document listing berthed vessels, vessels
at anchorage and expected arrivals, each with cargo, tonnage and an event time (berthing,
anchoring or estimated time of arrival). A scheduled job downloads it three times a day and
archives every document unchanged, so that every derived table can be rebuilt from the raw
files. The document header gives the time the document was \emph{generated}, not the time
its tables were updated. A change of situation is therefore detected on a fingerprint of
the table contents, never on the file or its header. The second public source is the port's
quarterly traffic statistics, which give tonnage landed by commodity.

\subsection{From observations to calls}
Each snapshot yields one observation per vessel. Observations are linked into calls by
vessel name, with a status sequence (expected, at anchorage, berthed) that may not move
backwards. A vessel that disappears briefly and reappears with the same cargo is attached
to the same call. Nothing is corrected silently: 13 anomaly codes flag revised arrival
estimates, berth changes, cargo changes and calls already in progress when the series
starts (left-censored).

A departure is never observed: a vessel is listed in one snapshot and absent from the next.
Each departure is therefore an interval between the last snapshot listing the vessel and the
first one that does not \cite{TODO-interval-censoring}. Durations computed on the lower end
of the interval are lower
bounds, and the width of the interval, set by the irregular timing of the scheduled job, is
published with each call. Waiting and berth times use the published event times, which are
actual movement times.

\subsection{What the series identifies, and what it does not}
The series identifies quantities measured \emph{at the port}. The apparent unloading rate,
published tonnage divided by berth occupancy, is 24\,000~t/day; it comes from 3 iron-ore
pellet calls, all at one berth and of one vessel size. Two waiting times at anchorage before
the ore berth are measured, 183.2~h and 159.4~h. Among 13 bulk carriers that waited at
anchorage, 6 of the 8 pairs that waited at the same time berthed in the reverse order of
arrival (75\%, Wilson interval [41; 93]\%) \cite{TODO-wilson-interval}. The test cannot speak
to the queue at a single berth: all 8 pairs berthed at different berths, and 5 of them carried
cargo of different categories.

The series does not identify quantities set \emph{before} the port. Expected arrivals are
announced about 1.5 days ahead, against an assumed mean supply lead time of 22 days. The
spread between announced and actual arrival gives a lead-time coefficient of variation of
0.028, a lower bound with no value as an estimate of the ordering lead time. The order of
first announcement against the order of arrival is reversed for 2 of 148 pairs (1\%, [0; 5]\%)
among 18 bulk carriers, which again measures arrival-estimate error, not overtaking between
orders. The ordering policy and the lead-time distribution remain unobserved.

From the quarterly statistics, the dispersion index of pellets landed per quarter is 0.83
(bootstrap interval [0.21; 1.73], 16 quarters) \cite{TODO-bootstrap}; it falls to 0.45 when
one anomalous quarter is removed. Read as the squared coefficient of variation of
inter-arrival times, it suggests near-Poisson arrivals, but that reading holds only if
arrivals are exogenous, which is precisely what Section~\ref{sec:trap} questions.

% =========================================================================================
\section{The trap: no steady state under demand-matched exogenous arrivals}
\label{sec:trap}

\subsection{Statement}
Consider a plant that consumes at rate $d$ while its feedstock inventory is positive and
stops otherwise. Vessels of i.i.d.\ sizes $Q_i$ arrive at the epochs of a renewal process
$N(t)$ whose mean rate is calibrated on demand, $\lambda = d/\mathbb{E}[Q]$, with
inter-arrival squared coefficient of variation $c_a^2$ and size squared coefficient of
variation $c_Q^2$. When the stockyard is full, a vessel keeps its cargo on board, at berth
or at anchorage, until space frees up: no supply is ever lost. Let $S_0$ be the initial
stock and
\begin{equation}
X(t) = S_0 + \sum_{i=1}^{N(t)} Q_i - d\,t
\end{equation}
the net input. The inventory is $I(t) = X(t) + L(t)$, where the lost production
\begin{equation}
L(t) = \max\Bigl(0,\; \sup_{u\le t} \bigl(-X(u)\bigr)\Bigr)
\end{equation}
is the one-sided regulator of $X$ at zero \cite{TODO-reflected-processes}. Define the
service level over horizon $T$ as $\beta(T) = 1 - L(T)/(dT)$, the fraction of target
production not lost for lack of feedstock.

\emph{Proposition.} Under these assumptions, $L(T)/\sqrt{T}$ converges in distribution to
$\sigma \sup_{u\le 1} B(u)$, where $B$ is a standard Brownian motion and
$\sigma^2 = d\,\mathbb{E}[Q]\,(c_a^2 + c_Q^2)$. Hence $1-\beta(T)$ decays as $T^{-1/2}$, the
cumulative deficit grows as $T^{1/2}$, $\beta(T)\to 1$ without reaching any limit below one,
and the stockyard capacity does not enter $L$.

\emph{Sketch.} The drift of $X$ is zero by calibration. By the functional central limit
theorem for renewal-reward processes \cite{TODO-fclt-renewal}, $(X(Tu)-S_0)/\sqrt{T}$
converges to $\sigma B(u)$. The regulator map is continuous, so $L(T)/\sqrt{T}$ converges to
the regulator of $\sigma B$ at time one, whose law is that of $\sigma\sup_{u\le1}B(u)$
\cite{TODO-reflection-principle}. The initial stock only shifts the barrier by a fixed
tonnage, whose weight in $\beta(T)$ vanishes as $1/T$. Since no supply is ever lost, the
upper bound of the inventory never acts on $X$, so the capacity cannot appear in $L$.
\hfill$\square$

Three consequences follow. Any service level read at a finite horizon is a transient. The
level read after one year mixes the corridor with the horizon. And storage capacity is
inert unless the model loses or diverts supply when it is full, which a model that lets
vessels wait does not.

\subsection{Measurement}
The simulated baseline, with Poisson arrivals, was replicated 20 times and each replication
was read at horizons from 1 to 50 years along the same trajectory. Fitting the mean deficit
as $aT^p$ gives $p = 0.578$ [0.434; 0.718] \emph{with} the nominal initial stock, and
$p = 0.492$ [0.367; 0.599] \emph{with zero} initial stock. Both intervals contain $1/2$ and
exclude 0 (bounded deficit) and 1 (stationary regime with losses); the initial stock steepens
the apparent slope because it shrinks the deficit at short horizons only, so the second
exponent is the one that describes the random walk itself. The service level rises from
0.858 at one year to 0.971 at fifty years without levelling off (Fig.~\ref{fig:thesis}, red).
Holding the initial stock fixed and multiplying both
storage capacities by 12 changes production by less than 0.002\% at horizons of 1, 4 and 12
years. An earlier result of the project, in which larger storage appeared to reduce losses
markedly, came entirely from the initial stock, which the scenarios derived from capacity.

\begin{figure}[t]
\centering
\includegraphics[width=\columnwidth]{\figdir/thesis.pdf}
\caption{Service level against the simulated horizon, same corridor, same 20 paired
replications per model, each read along one trajectory (bands: 95\% intervals). Red: exogenous
Poisson arrivals at the demand rate; $\beta$ keeps rising with no plateau. Blue: endogenous
$(s,Q)$ supply; $\beta$ stabilises, and the shaded strip is its limit $\beta_\infty$ estimated by
batch means over a single 400-year run.}
\label{fig:thesis}
\end{figure}

\begin{figure}[t]
\centering
\input{\figdir/structure.tex}
\caption{Structure of the corridor model. Black: material flow from sea to furnace. Red
dashed: the two blockings, which propagate upstream (a full stockyard holds the vessel at the
berth; a full plant stock holds the rakes). Blue: the supply loop of the endogenous model: the
inventory position, which counts everything afloat, at the port, on rail and at the plant, is
reviewed at each consumption, and a vessel is chartered when it falls below $s$. The exogenous
model has no loop: vessels come from an independent process at the demand rate (grey dotted).}
\label{fig:structure}
\end{figure}

% =========================================================================================
\section{The corrective model}
\label{sec:model}

\subsection{Endogenous supply}
The corrected model, written in SimPy \cite{TODO-simpy}, keeps the port, berth, rail and
plant processes of the original and replaces the arrival process by a continuous-review
$(s,Q)$ policy \cite{TODO-sQ-inventory-theory} (Fig.~\ref{fig:structure}). Whenever the inventory position falls below
the reorder point $s$, vessels are chartered one at a time ($Q$ = one vessel) until it is back
at or above $s$. The position counts the stockyard, the
plant stock, rail wagons in transit and everything afloat: vessels at sea, at anchorage, and
the cargo still on board at berth. Without the last term, the position would drop by a full
vessel at berthing and the policy would order for nothing. The policy is reviewed only when
the plant consumes, since unloading, rail transfer and arrivals move material without
removing it. Each order draws its vessel size, unloading rate and lead time from separate
random streams, so that the $n$-th order has the same three values in every configuration
compared, which makes paired comparisons tight.

All non-constant distributions are gamma laws with a given mean and coefficient of
variation. The lead time has a mean of 22 days [15; 30] and a coefficient of variation of
0.25 [0.10; 0.60]; the reorder point is 30 or 35 days of mean consumption, in a range of
[15; 60], set at the cost optimum of each variant. None of these three parameters has a
public source (Section~\ref{sec:limits}).

\subsection{Discrete outages and overtaking}
The plant has a fixed annual availability. In the \emph{continuous} variant (C) it runs at
that average rate; in the \emph{discrete} variant (D) it alternates between full rate and
unplanned outages with a mean duration of 7 days [2; 30], at the same availability. D is the
most irregular reading of the availability figure, an upper bound on consumption
irregularity. Lead times are drawn independently for each order, so a late vessel can be
overtaken by the next one. Two hypotheses are kept side by side: overtaking \emph{permitted}
(p), realistic for independent charters from different origins, and \emph{forbidden} (i), where
each vessel arrives no earlier than the previous one. The public series cannot separate them
(Section~\ref{sec:reconstruction}), so every result is reported under the four variants Cp,
Ci, Dp and Di.

\subsection{Evidence chain}
Table~\ref{tab:evidence} lists the checks. The berth queue alone was checked against the
exact Pollaczek--Khinchine mean wait \cite{TODO-pollaczek-khinchine}, in a scenario where the
storage is made large enough that service times are independent of the queue. The supply
model was then confronted with the quarterly dispersion of landed tonnage. With continuous
consumption it is refuted: it is far too regular, and the gap is not an artefact of trend or
vessel size. With discrete outages it is \emph{not refuted} on 2023--2025: 118 of 126
combinations of outage duration, storage level, lead-time variability and overtaking keep
the observed value inside the simulated 5--95\% band, and only the shortest outages (2 days)
are rejected. Three reservations bound this result. The test has low power: the simulated
bands span almost an order of magnitude (Fig.~\ref{fig:validation}). The reference period
2023--2025 was chosen \emph{after} seeing the role of 2022, which only outages of 30 days
reproduce (15 of 126 combinations, all at 30 days). And other mechanisms could produce a
similar dispersion. The test does not discriminate the lead-time variability either:
18, 19, 21, 21 and 21 of 21 combinations pass at coefficients of variation 0.10, 0.25, 0.40,
0.50 and 0.60.

\begin{table}[t]
\caption{Evidence chain for the corrective model}
\label{tab:evidence}
\centering
\footnotesize
\setlength{\tabcolsep}{3pt}
\begin{tabular}{@{}>{\raggedright\arraybackslash}p{2.9cm}>{\raggedright\arraybackslash}p{3.3cm}>{\raggedright\arraybackslash}p{1.6cm}@{}}
\toprule
Check & Statistic & Outcome \\
\midrule
Berth queue vs.\ Pollaczek--Khinchine (200 runs $\times$ 8 years) & mean wait 37.96~$\pm$~1.40~h vs.\ 38.53~h & consistent \\
Continuous consumption vs.\ traffic data (8 seeds $\times$ 100 years) & dispersion index, median 0.067--0.089 vs.\ 0.83; $P(I \ge 0.45)=0$ & refuted \\
Same, within-year CV$^2$ of tonnage & 0.174 observed vs.\ [0.006; 0.032] simulated & refuted \\
Discrete outages, 2023--2025 (4 seeds $\times$ 60 years) & observed value inside 5--95\% band & not refuted (118/126) \\
Discrete outages, 2022--2025 & same & 15/126, all at 30-day outages \\
\bottomrule
\end{tabular}
\end{table}

\begin{figure}[t]
\centering
\includegraphics[width=\columnwidth]{\figdir/validation.pdf}
\caption{Within-year squared coefficient of variation of quarterly tonnage landed, simulated
(median and 5--95\% band, discrete outages, lead-time coefficient of variation 0.25,
overtaking permitted, three storage levels) against the value observed on 2023--2025. The
observed value lies inside the bands except for the shortest outages, but the bands are
wide: the test has low power, and the period was chosen after the data were seen.}
\label{fig:validation}
\end{figure}

\begin{figure*}[t]
\centering
\includegraphics[width=\textwidth]{\figdir/factorial.pdf}
\caption{Effects of a second berth and of high storage on the corrected model, expansion
scenario, 200 paired replications of 16 years, four variants, 95\% intervals (often smaller
than the markers on the cost panels). Effects on the service level are in percentage points.
Costs are illustrative: the unit costs (margin per tonne, demurrage rate, pellet price,
carrying rate) have no local source, and demurrage is counted on all anchorage time without
laytime, hence an upper bound.}
\label{fig:factorial}
\end{figure*}

\subsection{Protocol}
Each run is warmed up for 4 years and measured over the next 16. The warm-up is handled by
differencing cumulative counters at the end of the warm-up rather than by resetting them,
which keeps both mass balances exact \cite{TODO-steady-state-simulation}. Configurations are
compared on paired replications sharing a master seed, with 95\% intervals on the paired
difference.

% =========================================================================================
\section{Results}
\label{sec:results}

\subsection{Stationarity recovered}
With endogenous supply, the service level converges. Read on the same corridor, with the
same protocol as the exogenous model, it stays between 0.987 and 0.991 at every horizon, and
no paired difference between successive horizons beyond 8 years is significant
(Fig.~\ref{fig:thesis}, blue). The deficit now grows linearly: fitted on horizons of 8 years
and more, its exponent is 1.048 [0.952; 1.134], the signature of a stationary regime with
losses. Batch means over a single 400-year run \cite{TODO-batch-means} give
$\beta_\infty = 0.9867 \pm 0.0012$ in the current configuration and $0.9912 \pm 0.0008$ in
the expansion scenario, and independent replications agree with them. The ranking of the two scenarios depends on overtaking: forbidding it lowers the
expansion scenario from 0.9917 to 0.9749, because with more orders in transit an overtaking
vessel more often fills a gap a late one would have opened \cite{TODO-order-crossover}. The
same feedback that restores a steady state also changes what dispersion of arrivals means:
it becomes an output of the policy and of the lead-time variability instead of an input.

\subsection{What the corrected model says about levers}
The following is an illustration of what the model implies, not a recommendation: the
storage capacities and the lead-time parameters are not sourced. Fig.~\ref{fig:factorial}
shows a $2\times2$ design (one or two ore berths, base or high storage) in the expansion
scenario, with 200 paired replications of 16 years in each of the four variants.

\emph{Storage does not serve production.} High storage changes the service level by
$-0.007$ to $-0.022$ points with one berth, nowhere a gain. This holds in all four variants.

\emph{The second berth acts on the queue, not on production.} Its effect on the service level
is significant but negligible, $+0.005$ to $+0.058$ points with base storage. It removes most
of the anchorage wait, and with it demurrage: $-1.96$ to $-3.56$~M~USD/year, and $-1.61$ to
$-3.03$~M~USD/year of total cost against a base-case total of 4.96 to 8.43~M~USD/year. These
amounts rest on unit costs without a local source and on demurrage counted as an upper bound;
they are gross of the investment and operation of a second berth, which is itself not an
observed configuration.

\emph{Stock-outs follow late vessels, not plant restarts.} Discrete outages lower the service
level by 1.3 to 2.1 points and multiply stock-out episodes: 0.15 and 0.85 episodes per year
in the continuous variants against 2.50 and 5.42 in the discrete ones. A natural explanation,
that the plant runs out after a restart because it then consumes at full rate, is rejected:
the relative risk of a stock-out in the 22 days after a restart is 0.20 and 0.30. A lead-time
delay beyond its 90th percentile precedes 85 to 98\% of episodes, against 21 to 23\% of
arbitrary instants. Counting the reorder point in days of full-rate rather than mean
consumption closes 100\% and 92\% of the continuous--discrete gap in episodes; in tonnes,
however, it is merely a higher reorder point, and it adds 89 and 82~kt of average stock.

\emph{The reorder point is the lever.} The cost-optimal reorder point grows with lead-time
variability, from 27.5 to 44.4 days (Ci) and from 31.7 to 44.8 days (Di) as the coefficient
of variation goes from 0.10 to 0.60, but much less when overtaking is permitted, from 27.5
to 35.2 days (Cp) and from 31.7 to 36.1 days (Dp). Propagating the parameter ranges through 500 draws, the long-run service level at
a fixed reorder point has a 5th percentile between 85.8\% and 97.2\% across variants; with
the reorder point adapted to the lead-time variability, between 92.3\% and 98.8\%. The
adapted figure assumes that the operator knows its lead-time variability, which this study
cannot measure.

% =========================================================================================
\section{Limitations and conclusion}
\label{sec:limits}

The analytic argument assumes no supply is ever lost and a finite-variance arrival
process. A model that diverts or rejects vessels when storage is full has an upper barrier,
and storage then matters; so does a model with any feedback from stock to arrivals.

The corridor study has several limitations. The two storage capacities are unknown and are
explored over wide ranges. The lead-time mean and variability and the reorder point have no
public source, and the snapshot series cannot measure them. The validation is a
non-refutation of low power on a period chosen after seeing the data, and discrete outages
are the most irregular reading of the plant availability. The continuous-review policy is
the most reactive possible, so its stock-outs are a lower bound when charters are in fact
fixed by monthly or quarterly programmes. The berth queue is first-come first-served, which
the observed series neither confirms nor refutes at a single berth. The unloading rate rests
on 3 calls. All monetary amounts use unit costs without a local source and an upper-bound
demurrage count. The overtaking result on the expansion scenario was measured by a
substitution script before the option existed in the model, and has not been recomputed.

The general lesson is simple and cheap to check. Before reading a service level from a
corridor model with exogenous arrivals, plot it against the simulated horizon: if it keeps
rising, there is nothing stationary to read, and the levers compared at that horizon are
compared on a transient. Giving the supply a feedback on the inventory position, even a
crude $(s,Q)$ rule, is enough to restore a steady state.

% =========================================================================================
\section*{Code and data availability}
The code, the scenario files, the archived public data and the scripts that produce every
number and figure of this paper are available at
\url{https://github.com/abdessamed-britah/djendjen-bellara-twin}; the results were frozen at
the tag \texttt{v0.2-resultats}. The raw port snapshots are archived unchanged as release
assets of that repository, and the traffic statistics as the published web pages. The
repository documentation is in French; \texttt{paper/NOMBRES.md} maps every number of this
paper to its line in the results table.

\section*{Acknowledgment}
The author used an AI assistant (Claude, Anthropic, through Claude Code) to write and test
the simulation and analysis code, to run the analyses and to draft this manuscript. All
numbers were produced by versioned scripts and are traceable to the repository. The author
reviewed the content and takes full responsibility for it.

\bibliographystyle{IEEEtran}
\bibliography{refs}

\end{document}
